\documentclass{ximera}

\input{../../preamble.tex}

\title{Quadratic  Behavior}

\begin{document}

\begin{abstract}
%Stuff can go here later if we want!
\end{abstract}
\maketitle


 
 




We have two different viewpoints when examining functions.


$\blacktriangleright$ \textbf{Numbers:}  Functions consist of individual pairs of individual numbers.  When we \textbf{\textcolor{purple!85!blue}{evaluate}} functions, we think of each pair one at a time.  We calculate with formulas, one number at a time. The graphs of functions are individual dots plotted to represent points, which visually encode each function pair. 







$\blacktriangleright$ \textbf{Intervals:} As we compare function values, we notice patterns, which we call \textbf{\textcolor{purple!85!blue}{behavior}}. We naturally group many pairs together where the function has similar behavior.  Collecting such pairs forms domain intervals and we begin thinking in terms of intervals.





In particular, we calculate \textbf{\textcolor{blue!55!black}{rate of change}} over intervals, which measures behavior. 



The rate of change of the function $f$ over the interval $[a, b]$ is given by $\frac{f(b) - f(a)}{b - a}$.




However, intervals are just collections of individual numbers.  Is there a way to translate this ``rate of change over an interval'' idea down to individual numbers?

The answer is yes.  









\subsection*{Making Sense}

We will use limits to make sense out of 


\[
\frac{f(a) - f(a)}{a-a}
\]




Limits give us expected values. We can use limits to give us an expected value for $\frac{f(a) - f(a)}{a-a}$.



Our plan is to calculate $\frac{f(b) - f(c)}{b-c}$ for intervals continaing $a$, $a \in (b,c)$.

We'll let this interval shrink down closer and closer to $a$.

As the interval shrinks, we'll watch the values of $\frac{f(b) - f(c)}{b-c}$.

From this we can determine an ``expected value''.


We will call this value \textbf{the derivative of f at a : $f'(a)$}






















We will begin this thought experiment here with quadratic functions. We'll then extend this thought a bit throughout the course. Then Calculus will answer the question fully with the derivative.















\subsection*{Learning Outcomes}

\begin{sectionOutcomes}
In this section, students will 

\begin{itemize}
\item develop the $iRoC$ function (the derivative, $f'$).
\item analyze quadratic functions.
\end{itemize}
\end{sectionOutcomes}










\begin{center}
\textbf{\textcolor{green!50!black}{ooooo-=-=-=-ooOoo-=-=-=-ooooo}} \\

more examples can be found by following this link\\ \link[More Examples of Quadratic Behavior]{https://ximera.osu.edu/csccmathematics/precalculus1/precalculus1/quadraticBehavior/examples/exampleList}

\end{center}





\end{document}
